\section{复算与人工智能使用说明}
支撑材料给出三问的完整调度文本。问题一只含原始节点顺序；问题二、三分别提供扩展节点顺序、初始地址偏移与按换出顺序列出的新偏移。新增节点每对由换出及换入组成，重新排列后也保持列表与节点编号一致。复算应从原图和这些文本重新建立依赖、空间状态及执行时间，不能只读取汇总数值。

程序以 Python 实现。图读取、顺序生成、地址分配、流水重排与附件重放分模块组织；随机种子不适用，因为候选规则及平局处理均为确定性。候选保留失败记录，论文表中的最终值来自同一组六图方案。本研究不把程序正常结束或结果文件存在当作数值正确的充分条件。

模型推导、程序及文字分析使用 OpenAI Codex 辅助。工具精确模型标识和版本发布日期未核实，参赛队员独立复核尚未完成。本阅读稿不构成人工原创、独立审查或正式提交资格的声明。若用于参赛，应由队员依据适用规定理解、核验并完成真实工具记录与使用声明\cite{ai}。
